\section{问题一模型建立及求解}

% v8.0.1 A196-inspired question chapter.
% 复杂问题默认形成“模型建立 → 模型求解 → 求解结果 → 结果的分析与验证”的局部闭环；
% 这些是功能阶段，不是不可修改的标题。若更专业的题目专属标题能更准确暴露数学对象，可直接替换。
% 简单解析题或直接计算题可以合并、删减非必要阶段。

\subsection{模型建立}
% 建议顺序：对象与当前缺口 → 变量/状态 → 机理或定义 → 目标/判定关系 → 约束 → 可计算模型。
% 优化类问题可按需设置三级标题：决策变量、目标函数、约束条件、模型汇总；
% 机理类问题则应按真实对象设置专业标题，不机械套“决策变量”。

设……，由前述共享关系/本问条件可得
\begin{equation}
    F(x,\theta)=0,
    \label{eq:q1-core-relation}
\end{equation}
式~\eqref{eq:q1-core-relation} 给出……，并作为后续求解的直接输入。

% 命题只在承担等价变换、单调性、可行性、降维或边界判定等真实作用时使用。
% \begin{hskproposition}{……}
% \label{prop:q1-example}
% ……
% \begin{hskproof}
% ……
% \end{hskproof}
% \end{hskproposition}

% 对复杂优化模型，模型建立末尾可进行核心模型收束。
% 目标函数必须单独展示，不进入约束大括号：
\begin{equation}
    \min_{\mathbf{x}} f(\mathbf{x})
    \label{eq:q1-objective}
\end{equation}

\begin{equation}
\text{s.t.}\quad
\left\{
\begin{aligned}
    g_i(\mathbf{x}) &\le 0, \quad i=1,\ldots,m, \\
    h_j(\mathbf{x}) &= 0, \quad j=1,\ldots,n, \\
    \mathbf{x} &\in \Omega .
\end{aligned}
\right.
\label{eq:q1-constraints}
\end{equation}
% 非优化题必须删除上述 optimization 示例，改成真实状态方程、概率映射、边界条件或解析关系。
% “模型汇总”是模型建立末尾的收束动作，不默认单列“核心模型汇总”二级标题。

\subsection{模型求解}
% 从“最终模型已经化成什么计算问题”开始，而不是先介绍算法百科。
% 写清：剩余计算困难 → solver/离散/搜索/分解方法为什么适配 → 关键参数、初值、精度、终止条件 → 输出映射。
% 若循环、分支、修复、接受/拒绝或终止逻辑本身是关键证据，可在解释输入与原因后再给算法步骤/伪代码。

\subsection{求解结果}
% 参考优秀论文的局部证据链：主结果立即在本问展示，图/表之后紧跟决定性数值、趋势和机制解释。
% 不要连续堆多张图后再统一解释；本节末尾应自然、直接回答问题一。

\subsection{结果的分析与验证}
% 当存在真实风险需要检查时启用：离散精度、参数敏感性、鲁棒性、多初值/多 seed、边界扰动、
% 替代算法/模型、残差/误差区间、外样本等。每项验证说明：检验什么风险 → 改变什么 → 指标 → 结果变化 → 支持边界。
% 若问题一属于简单解析/直接计算且没有实质验证必要，可删除整个小节。

% 按实际小问复制为 07_question2.tex、08_question3.tex 等。
% 后问继承前文已经定义的共享模型，只写新增变量、目标、约束、算法和证据，不从头复制模型准备。
